\Needspace{18\baselineskip}
\section{Continuous representatives and pullback}
\label{sec:continuous-bc-representatives}

A finite divisorial sum is pulled back through the local equations of its divisors. For an infinite sum we need pullback to be continuous on closed positive currents whose Bott--Chern classes vary. Solving a scalar equation for the Gauduchon trace selects a smooth representative in each class; the complementary normalized potentials then converge in $L^1$, and \cref{lem:modification-lone-pullback} shows that their pullbacks converge as well.

In this section, $Y$ denotes an arbitrary connected compact complex manifold of dimension $n\geq1$, and $\gamma$ is a Gauduchon metric on $Y$; in dimension one, any Hermitian metric will do. Fix smooth norms on tensor fields induced by $\gamma$. Put $V_\gamma=\int_Y\gamma^n$ and, for a real $(1,1)$-form or current $B$, define its \emph{scalar trace} $\sigma_\gamma(B)$ by
\[
 \sigma_\gamma(B)\,\gamma^n=B\wedge\gamma^{n-1}.
\]
For a current this is an identity of distributions. As in the proof of \cref{thm:gauduchon}, we use the classical elliptic Fredholm theorem, elliptic regularity, and the continuity on distributions of the normalized elliptic inverse.

\Needspace{18\baselineskip}
\begin{lemma}\label{lem:constant-trace-bc-representative}\index{Bott--Chern class!constant-trace representative}
Every closed real $(1,1)$-current $T$ has a unique decomposition
\[
 T=B_\gamma(T)+\ddc v_\gamma(T),\qquad
 \int_Yv_\gamma(T)\,\gamma^n=0,
\]
where $v_\gamma(T)$ is a real distribution and $B_\gamma(T)$ is a smooth closed real form with constant scalar trace. Both assignments are linear and continuous from weak convergence of currents to convergence of distributions. The forms $B_\gamma(T)$ lie in a finite-dimensional space, and their convergence takes place in $C^\infty$.

This space maps isomorphically onto $H^{1,1}_{\mathrm{BC}}(Y,\mathbb R)$. If $T\geq0$, then $v_\gamma(T)$ is quasi-plurisubharmonic. On the cone of positive closed currents, $v_\gamma(T)$ depends continuously on $T$ in $L^1$, and for every nonnegative integer $l$,
\begin{equation}\label{eq:constant-trace-representative-bound}
 \|B_\gamma(T)\|_{C^l}+\|v_\gamma(T)\|_{L^1(\gamma^n)}
       \leq C_l\int_YT\wedge\gamma^{n-1},
\end{equation}
where the constants $C_l$ depend only on the fixed metric and norms.
\end{lemma}
\begin{proof}
Set $L_\gamma f=\sigma_\gamma(\ddc f)$. We use the inverse $G_\gamma$ of \eqref{eq:gaud-trace-green}, normalized with the volume form $\gamma^n/V_\gamma$. It acts on top-degree distributions of total mass zero and is continuous by the elliptic estimates for the operator and its transpose.

\emph{Step 1: construction.} Define
\[
 c_T=\frac{\int_YT\wedge\gamma^{n-1}}{V_\gamma},\qquad
 v_\gamma(T)=G_\gamma\bigl((\sigma_\gamma(T)-c_T)\gamma^n\bigr),\qquad
 B_\gamma(T)=T-\ddc v_\gamma(T).
\]
The current $B_\gamma(T)$ is closed and has scalar trace $c_T$. We now prove its smoothness and the finite-dimensionality assertion without using any K\"ahler identity.

\Needspace{7\baselineskip}
\emph{Step 2: an elliptic system.} Let $J$ be the complex structure acting on real covectors. On real $(1,1)$-forms consider the first-order operator
\[
 \mathscr P B=(\mathrm{d}B,\mathrm{d}\sigma_\gamma(B)).
\]
Its principal symbol at a nonzero real covector $\xi$ is $B\mapsto(\xi\wedge B,\xi\,\sigma_\gamma(B))$, and we claim that it is injective. If $\xi\wedge B=0$, then the $(1,1)$ condition forces $B$ to be a real multiple of $\xi\wedge J\xi$. To check this, choose unitary real coordinates with $\xi=\mathrm{d}x_1$: every term of $B$ contains $\mathrm{d}x_1$, and invariance under $J$ forces every term to contain $\mathrm{d}y_1$ as well. If moreover $\xi\,\sigma_\gamma(B)=0$, then $\sigma_\gamma(B)=0$; since the scalar trace of $\xi\wedge J\xi$ is nonzero, the multiple is zero.

Let $\mathscr P^*$ be the formal adjoint of $\mathscr P$ for the fixed metric and volume. Since the symbol of $\mathscr P$ is injective, $\mathscr P^*\mathscr P$ is elliptic. By elliptic regularity and compactness, its distributional kernel is a finite-dimensional space of smooth forms. Since $\mathscr P B_\gamma(T)=0$, the form $B_\gamma(T)$ belongs to this space. Conversely, the kernel consists exactly of the closed forms with constant scalar trace: integrating $\langle\mathscr P^*\mathscr P B,B\rangle$ gives $\mathscr P B=0$, and connectedness makes the trace constant. Denote this kernel by $\mathscr H_\gamma$.

\emph{Step 3: Bott--Chern representatives.} The construction is linear and continuous on distributions. Convergence of the components $B_\gamma(T)$ therefore takes place in the finite-dimensional space $\mathscr H_\gamma$, hence in every smooth norm. If $T=\ddc f$ for a distribution $f$, then $c_T=0$ and the inverse gives $v_\gamma(T)=f-V_\gamma^{-1}\int_Yf\gamma^n$, so $B_\gamma(T)=0$. Applied to smooth closed forms, this shows that $\mathscr H_\gamma$ contains exactly one representative of each Bott--Chern class. Uniqueness of the decomposition can also be seen directly: if two decompositions are given, $L_\gamma$ of the difference of the two potentials is a constant, integration forces this constant to be zero, and the normalized kernel of $L_\gamma$ is zero.

\emph{Step 4: positive currents.} Suppose now that $T\geq0$. The inequality $\ddc v_\gamma(T)\geq-B_\gamma(T)$ gives a quasi-plurisubharmonic representative of the distribution $v_\gamma(T)$. Let positive closed currents $T_j$ converge weakly to $T$. Their components $B_\gamma(T_j)$ converge smoothly, so the potentials $v_\gamma(T_j)$ have a common smooth lower curvature bound. By \cref{thm:normalized-compactness}, after subtracting their suprema they form an $L^1$-precompact family. If $a_j=\sup_Yv_\gamma(T_j)$, the zero-mean normalization gives
\[
 0\leq a_j V_\gamma=\int_Y\bigl(a_j-v_\gamma(T_j)\bigr)\,\gamma^n,
\]
and the right-hand side is uniformly bounded. Hence the potentials $v_\gamma(T_j)$ themselves are $L^1$-precompact. Every cluster point equals their distributional limit $v_\gamma(T)$, so the whole sequence converges in $L^1$.

Finally, the positive closed currents with trace mass $\int_YT\wedge\gamma^{n-1}=1$ form a weakly compact set: positivity bounds their coefficient measures, weak limits remain positive and closed, and testing with $\gamma^{n-1}$ preserves the trace mass. The smooth and $L^1$ continuity just proved bounds both norms on this set, and positive homogeneity gives \eqref{eq:constant-trace-representative-bound}. A positive current of trace mass zero is zero, so this case is included.
\end{proof}

The smooth representatives give a common curvature bound for converging normalized potentials. This is exactly the hypothesis needed to pull back their $L^1$ convergence.

\Needspace{13\baselineskip}
\begin{proposition}\label{prop:closed-current-modification-continuity}\index{Closed positive current!pullback continuity}
Let $g:Z\to Y$ be a smooth compact modification, and fix a Gauduchon metric $\gamma'$ on $Z$. Pullback of positive closed $(1,1)$-currents, defined through their local potentials, is weakly continuous. Moreover, there is a constant $C_g$, depending only on the modification and the fixed metrics, such that
\begin{equation}\label{eq:modification-linear-mass-bound}
 \int_Zg^*T\wedge(\gamma')^{n-1}
       \leq C_g\int_YT\wedge\gamma^{n-1}
       \qquad(T\geq0,\ \mathrm{d}T=0).
\end{equation}
\end{proposition}
\begin{proof}
Write $T=B_\gamma(T)+\ddc v_\gamma(T)$. The pullback is
\[
 g^*T=g^*B_\gamma(T)+\ddc(v_\gamma(T)\circ g).
\]
This current is positive and closed, since it is locally $\ddc$ of plurisubharmonic functions, and the composed potential is not identically $-\infty$ because $g$ is surjective. The formula agrees with the definition through local potentials, so it does not depend on the decomposition.

If $T_j\to T$, the smooth parts converge in $C^\infty$, and the potentials converge in $L^1(Y)$ with a common lower curvature bound. By \cref{lem:modification-lone-pullback}, their compositions with $g$ converge in $L^1(Z)$, and taking $\ddc$ proves the weak convergence of the pullbacks.

\Needspace{7\baselineskip}
Integration by parts on $Z$ gives
\[
 \int_Zg^*T\wedge(\gamma')^{n-1}
   =\int_Zg^*B_\gamma(T)\wedge(\gamma')^{n-1}.
\]
The absolute value of the right-hand side is at most a fixed constant times $\|B_\gamma(T)\|_{C^0}$. Since the left-hand side is nonnegative, \eqref{eq:constant-trace-representative-bound} gives \eqref{eq:modification-linear-mass-bound}.
\end{proof}

\Needspace{12\baselineskip}
\begin{corollary}\label{cor:countable-divisorial-pullback}\index{Divisorial current!countable pullback}
Let $g:Z\to Y$, $\gamma$ and $\gamma'$ be as in \cref{prop:closed-current-modification-continuity}. Let $D_j$ be prime divisors on $Y$, with integration currents $[D_j]$, and let $c_j>0$. If $S=\sum_j c_j[D_j]$ is a positive closed current of finite mass, then
\[
 g^*S=\sum_j c_j[g^*D_j]
\]
as positive currents of finite mass. The trace masses on $Z$ of the positive tails of this series tend to zero. In particular, $g^*S$ is again purely divisorial.
\end{corollary}
\begin{proof}
Apply the proposition to the increasing finite partial sums $S_N$ and to their positive closed tails $S-S_N$. Pullback is additive on positive closed currents, by addition of local potentials. Thus $g^*S-g^*S_N=g^*(S-S_N)\geq0$, and its trace mass is at most $C_g$ times the trace mass of $S-S_N$, which tends to zero. The pullback of a finite divisor is the usual pullback of its local equations, which identifies $g^*S_N$ with the corresponding partial sum of the series, and hence $g^*S$ with its limit. Its terms are integration currents of prime divisors with nonnegative coefficients, and regrouping this countable positive sum of finite mass does not change it.
\end{proof}
